\documentclass[12pt]{article}
\usepackage{amsmath, amssymb, amsthm}
\usepackage{centernot}
\usepackage{geometry}
\geometry{margin=2.5cm}

% Simple manual environments - no amsthm numbering involved
\newenvironment{satz}[1]{%
  \par\vspace{6pt}\noindent\textbf{Theorem #1.}\itshape\space
}{%
  \par\vspace{6pt}
}

\newenvironment{definition}[1]{%
  \par\vspace{6pt}\noindent\textbf{Definition #1.}\itshape\space
}{%
  \par\vspace{6pt}
}

\newenvironment{bemerkung}[1]{%
  \par\vspace{6pt}\noindent\textbf{Remark #1.}\itshape\space
}{%
  \par\vspace{6pt}
}

\renewcommand{\qedsymbol}{$\blacksquare$}
\setcounter{secnumdepth}{0}

\title{7 Single Variable Integration}
\date{}

\begin{document}

\maketitle

In this chapter the (Riemann) integral for functions of one (real) variable is introduced.
This is done first for step functions and then, with the help of a general theorem on the
extension of linear operators, for jump-continuous functions.

Some properties of the integral and important integration rules are derived.
The notion of integral is then extended to so-called improper integrals. As an important
example of improper integrals, the Euler gamma function is studied.

\section{7.1 The (Riemann) Integral}

In the following, let $(X,\|\cdot\|_X)$ always be a $\mathbb K$-Banach space and $I=[a,b]$ a compact perfect
interval, i.e.\ $a<b$.

\begin{definition}{7.1.1}
Let $f\in T(I,X)$ and $Z=(\alpha_0,\ldots,\alpha_n)$ an associated partition. Then
the \emph{integral} of $f$ with respect to $Z$ is defined by:
\[
\int_{(Z)} f
:=\sum_{i=1}^{n}(\alpha_i-\alpha_{i-1})\cdot f\!\left(\frac{\alpha_i+\alpha_{i-1}}{2}\right).
\]
\end{definition}

\par\vspace{6pt}\noindent\textbf{Lemma 7.1.2.}\itshape\space
Let $f\in T(I,X)$.
\begin{enumerate}
    \item There exists a minimal associated partition $Z_0$ of $I$.
    \item Every associated partition $Z_1$ of $I$ is a refinement of $Z_0$.
    \item $\int_{(Z)} f$ is independent of the particular associated partition $Z$.
\end{enumerate}
\upshape

\begin{proof}
\textbf{Ad 1,2:} Clear.

\medskip
\textbf{Ad 3:} Let $Z'$ be an elementary refinement of $Z=(\alpha_0,\ldots,\alpha_n)$, i.e.
\[
Z'=(\alpha_0,\ldots,\alpha_k,\gamma,\alpha_{k+1},\ldots,\alpha_n).
\]
Then:
\begin{align*}
\int_{(Z)} f
&=\sum_{i=1}^{k}(\alpha_i-\alpha_{i-1})f\!\left(\frac{\alpha_i+\alpha_{i-1}}{2}\right)
+(\alpha_{k+1}-\gamma+\gamma-\alpha_k)f\!\left(\frac{\alpha_k+\alpha_{k+1}}{2}\right)\\
&\quad+\sum_{i=k+2}^{n}(\alpha_i-\alpha_{i-1})f\!\left(\frac{\alpha_{i-1}+\alpha_i}{2}\right)\\[6pt]
&=\sum_{i=1}^{k}(\alpha_i-\alpha_{i-1})f\!\left(\frac{\alpha_{i-1}+\alpha_i}{2}\right)
+(\gamma-\alpha_k)f\!\left(\frac{\alpha_k+\gamma}{2}\right)
+(\alpha_{k+1}-\gamma)f\!\left(\frac{\gamma+\alpha_{k+1}}{2}\right)\\
&\quad+\sum_{i=k+2}^{n}(\alpha_i-\alpha_{i-1})f\!\left(\frac{\alpha_{i-1}+\alpha_i}{2}\right)\\[6pt]
&=\int_{(Z')} f.
\end{align*}
Since every associated partition $Z$ can be written as $Z=Z_0\cup Z_1\cup\cdots\cup Z_l$,
where $Z_0$ is the minimal associated partition and $Z_i$, $1\le i\le l$, are elementary
refinements, it follows by induction that
\[
\int_{(Z)} f=\int_{(Z_0)} f.
\]
\end{proof}

\begin{definition}{7.1.3}
Let $f\in T(I,X)$. Then the \emph{Riemann integral} of $f$ is defined by
\[
\int_I f:=\int_a^b f:=\int_{(Z)} f,
\]
where $Z$ is any admissible partition of $I$ for $f$.
\end{definition}

\par\vspace{6pt}\noindent\textbf{Lemma 7.1.4.}\itshape\space
The map $\int_a^b:T(I,X)\to X$, $f\mapsto\int_a^b f$ is linear, i.e.
\[
\int_a^b(f+g)=\int_a^b f+\int_a^b g,
\qquad
\int_a^b \alpha f=\alpha\int_a^b f
\qquad
\forall f,g\in T(I,X),\ \alpha\in\mathbb K.
\]
Moreover:
\[
\left\|\int_a^b f\right\|_X\le(b-a)\|f\|_\infty
\qquad
\forall f\in T(I,X).
\]
\upshape

\begin{proof}
The linearity of the integral follows directly from the definition.
Let $f\in T(I,X)$ and $Z=(\alpha_0,\ldots,\alpha_n)$ an associated partition of $I$. Then:
\[
\left\|\int_a^b f\right\|_X
=\left\|\int_{(Z)} f\right\|_X
=\left\|\sum_{i=1}^{n}(\alpha_i-\alpha_{i-1})f\!\left(\frac{\alpha_{i-1}+\alpha_i}{2}\right)\right\|_X
\le\sum_{i=1}^{n}(\alpha_i-\alpha_{i-1})\|f\|_{\infty,]\alpha_{i-1},\alpha_i[}
\le(b-a)\|f\|_\infty.
\]
\end{proof}

\begin{definition}{7.1.5}
Let $(Y,\|\cdot\|_Y)$ and $(Z,\|\cdot\|_Z)$ be normed $\mathbb K$-vector spaces.
\begin{enumerate}
    \item A linear map $A:Y\to Z$ is called \emph{bounded} if there exists $\alpha\in\mathbb R_{\ge 0}$
    such that
    \[
    \|Ax\|_Z\le\alpha\|x\|_Y
    \qquad\forall x\in Y.
    \]

    \item $\mathcal L(Y,Z):=\{A:Y\to Z: A\text{ is bounded}\}$ is a vector space.

    \item $\|A\|:=\|A\|_{\mathcal L(Y,Z)}:=\inf\{\alpha\in\mathbb R_{\ge 0}:\|Ax\|_Z\le\alpha\|x\|_Y,\ \forall x\in Y\}$
    is a norm on $\mathcal L(Y,Z)$.
\end{enumerate}
\end{definition}

\begin{bemerkung}{7.1.6}
\begin{enumerate}
    \item Let $A\in\mathcal L(Y,Z)$. Then
    \[
    \|A\|_{\mathcal L(Y,Z)}
    =\sup_{\substack{x\in Y\setminus\{0\}}}\frac{\|Ax\|_Z}{\|x\|_Y}
    =\sup_{\substack{x\in Y\setminus\{0\}\\\|x\|_Y\le 1}}\|Ax\|_Z
    =\sup_{\substack{x\in Y\\\|x\|_Y=1}}\|Ax\|_Z.
    \]

    \item Let $A\in\mathcal L(Y,Z)$. Then:
    \[
    \|Ax\|_Z\le\|A\|_{\mathcal L(Y,Z)}\cdot\|x\|_Y
    \qquad\forall x\in Y.
    \]

    \item Let $A:Y\to Z$ be linear. Then:
    \[
    A\text{ is uniformly continuous}
    \iff A\text{ is continuous}
    \iff A\text{ is continuous at }0
    \iff A\in\mathcal L(Y,Z).
    \]
\end{enumerate}
\end{bemerkung}

\begin{proof}
\begin{enumerate}
    \item We have:
    \begin{align*}
    \|A\|_{\mathcal L(Y,Z)}
    &=\inf\{\alpha\in\mathbb R_{\ge 0}:\|Ax\|_Z\le\alpha\|x\|_Y,\ \forall x\in Y\}\\
    &=\inf\{\alpha\in\mathbb R_{\ge 0}:\|Ax\|_Z\le\alpha\|x\|_Y,\ \forall x\in Y\setminus\{0\}\}\\
    &=\inf\left\{\alpha\in\mathbb R_{\ge 0}:\frac{\|Ax\|_Z}{\|x\|_Y}\le\alpha,\ \forall x\in Y\setminus\{0\}\right\}\\
    &=\sup_{\substack{x\in Y\setminus\{0\}}}\frac{\|Ax\|_Z}{\|x\|_Y}
    =\sup_{\substack{x\in Y\setminus\{0\}}}\left\|A\frac{x}{\|x\|_Y}\right\|_Z
    =\sup_{\substack{x\in Y\setminus\{0\}\\\|x\|_Y=1}}\|Ax\|_Z\\
    &\le\sup_{\substack{x\in Y\setminus\{0\}\\\|x\|_Y\le 1}}\|Ax\|_Z.
    \end{align*}
    On the other hand, let $x\in Y$ with $0<\|x\|_Y\le 1$. Then:
    \[
    \|Ax\|_Z
    \le\frac{1}{\|x\|_Y}\|Ax\|_Z
    =\left\|A\frac{x}{\|x\|_Y}\right\|_Z.
    \]
    Hence
    \[
    \sup_{\substack{x\in Y\setminus\{0\}\\\|x\|_Y\le 1}}\|Ax\|_Z
    \le\sup_{\substack{x\in Y\setminus\{0\}\\\|x\|_Y=1}}\|Ax\|_Z.
    \]
    The claim follows.

    \item For $x\neq 0$ this follows from
    \[
    \|Ax\|_Z
    =\frac{\|Ax\|_Z}{\|x\|_Y}\|x\|_Y
    \le\left(\sup_{\substack{w\in Y\setminus\{0\}}}\frac{\|Aw\|_Z}{\|w\|_Y}\right)\|x\|_Y.
    \]
    For $x=0$ the statement is trivial.

    \item $A$ uniformly continuous $\Longrightarrow$ $A$ continuous $\Longrightarrow$ $A$ continuous at $0$

    $\Longrightarrow$ $\exists\,\delta>0$ $\forall x\in Y$ with $\|x\|_Y\le\delta$: $\|Ax\|_Z\le 1$

    $\Longrightarrow$ $\exists\,\delta>0$ $\forall x\in Y$ with $\|x\|_Y=1$:
    $\|A\delta x\|_Z=\delta\|Ax\|_Z\le 1$, i.e.\ $\|Ax\|_Z\le\dfrac{1}{\delta}$

    $\Longrightarrow$ $A$ is bounded and $A\in\mathcal L(Y,Z)$

    $\Longrightarrow$ $\|Ax-Ay\|_Z=\|A(x-y)\|_Z\le\|A\|_{\mathcal L(Y,Z)}\cdot\|x-y\|_Y
    \le\dfrac{1}{\delta}\|x-y\|_Y$.
\end{enumerate}
\end{proof}

\par\vspace{6pt}\noindent\textbf{Lemma 7.1.7.}\itshape\space
Let $(Y,\|\cdot\|_Y)$ be a normed $\mathbb K$-vector space and $U\subset Y$ a subspace.
Then $\overline{U}$ is also a subspace.
\upshape

\begin{proof}
Exercise.
\end{proof}

\begin{satz}{7.1.8 (Extension of continuous linear operators)}
Let $(Y,\|\cdot\|_Y)$ be a normed $\mathbb K$-vector space, $U\subset Y$ a subspace,
$(Z,\|\cdot\|_Z)$ a $\mathbb K$-Banach space, and $A\in\mathcal L(U,Z)$.
Then there exists exactly one $\overline{A}\in\mathcal L(\overline{U},Z)$ with:
\begin{enumerate}
    \item $\overline{A}y=Ay\qquad\forall y\in U$,
    \item $\|\overline{A}\|_{\mathcal L(\overline{U},Z)}=\|A\|_{\mathcal L(U,Z)}$,
    \item $\overline{A}y=\lim_{n\to\infty}Ay_n\qquad\forall (y_n)_{n\in\mathbb N_0}\subset U:\ y=\lim_{n\to\infty}y_n$.
\end{enumerate}
$\overline{A}$ is called the \emph{extension} of $A$.
\end{satz}

\begin{proof}
\textbf{Uniqueness:}
Let $B,C\in\mathcal L(\overline{U},Z)$ be two extensions satisfying (1). Then their continuity implies
that for every $y\in\overline{U}$ with $y=\lim_{n\to\infty}y_n$ and $(y_n)_{n\in\mathbb N_0}\subset U$:
\[
By=B\!\left(\lim_{n\to\infty}y_n\right)=\lim_{n\to\infty}(By_n)=\lim_{n\to\infty}(Cy_n)=C\!\left(\lim_{n\to\infty}y_n\right)=Cy.
\]
Hence $B=C$.

\medskip
\textbf{Existence:}
Let $y\in\overline{U}$ and $(y_n)_{n\in\mathbb N_0}\subset U$ with $y=\lim_{n\to\infty}y_n$.
Then $(y_n)_{n\in\mathbb N_0}$ is a Cauchy sequence in $Y$ and
\[
\underbrace{\|Ay_n-Ay_m\|_Z}_{\in Z}
=\|A\underbrace{(y_n-y_m)}_{\in U}\|_Z
\le\|A\|_{\mathcal L(U,Z)}\|y_n-y_m\|_Y.
\]
Hence $(Ay_n)_{n\in\mathbb N_0}$ is a Cauchy sequence and, since $Z$ is a Banach space, it converges
(with limit $z:=\lim_{n\in\mathbb N_0}Ay_n$). Define $\overline{A}y:=z$.

\medskip
\textit{Independence of the chosen sequence:}
Let $(\tilde y_n)_{n\in\mathbb N_0}\subset U$ be another sequence in $U$ with $y=\lim_{n\to\infty}\tilde y_n$.
Let $\tilde z:=\lim_{n\to\infty}A\tilde y_n$. Then:
\[
\|z-\tilde z\|_Z
\le\|z-Ay_n\|_Z
+\underbrace{\|Ay_n-A\tilde y_n\|_Z}_{=\|A(y_n-\tilde y_n)\|_Z\le\|A\|\|y_n-\tilde y_n\|_Y\to 0}
+\|A\tilde y_n-\tilde z\|_Z
\xrightarrow[n\to\infty]{}0.
\]
Hence $z=\tilde z$.

\medskip
\textit{Linearity of\/ $\overline{A}$:}
Let $u,v\in\overline{U}$ and $(u_n)_{n\in\mathbb N_0},(v_n)_{n\in\mathbb N_0}\subset U$ with $u=\lim_{n\to\infty}u_n$,
$v=\lim_{n\to\infty}v_n$. Then:
\begin{align*}
\overline{A}(u+v)
&=\overline{A}\!\left(\lim_{n\to\infty}u_n+\lim_{n\to\infty}v_n\right)
=\overline{A}\!\left(\lim_{n\to\infty}(u_n+v_n)\right)
=\lim_{n\to\infty}\overline{A}(u_n+v_n)\\
&=\lim_{n\to\infty}A(u_n+v_n)
=\lim_{n\to\infty}(Au_n+Av_n)
=\lim_{n\to\infty}Au_n+\lim_{n\to\infty}Av_n \tag{7.1.1}\\
&=\overline{A}\!\left(\lim_{n\to\infty}u_n\right)+\overline{A}\!\left(\lim_{n\to\infty}v_n\right)
=\overline{A}u+\overline{A}v.
\end{align*}
Analogously one shows $\overline{A}(\alpha u)=\alpha\overline{A}u$.

\medskip
\textit{Norm preservation:}
Let $y\in\overline{U}$ with $y=\lim_{n\to\infty}y_n$, $(y_n)_{n\in\mathbb N_0}\subset U$. Then:
\[
\|\overline{A}y\|_Z
=\left\|\lim_{n\to\infty}Ay_n\right\|_Z
=\lim_{n\to\infty}\|Ay_n\|_Z
\le\lim_{n\to\infty}\|A\|_{\mathcal L(U,Z)}\|y_n\|_Y
=\|A\|_{\mathcal L(U,Z)}\|y\|_Y.
\]
Hence:
\[
\|\overline{A}\|_{\mathcal L(\overline{U},Z)}
=\sup_{\substack{y\in\overline{U}\\\|y\|_Y=1}}\|\overline{A}y\|_Z
\le\|A\|_{\mathcal L(U,Z)}.
\]
By (1):
\[
\|\overline{A}\|_{\mathcal L(\overline{U},Z)}
=\sup_{\substack{y\in\overline{U}\\\|y\|_Y=1}}\|\overline{A}y\|_Z
\ge\sup_{\substack{y\in U\\\|y\|_Y=1}}\|\overline{A}y\|_Z
=\sup_{\substack{y\in U\\\|y\|_Y=1}}\|Ay\|_Z
=\|A\|_{\mathcal L(U,Z)}.
\]
Together this gives
\[
\|\overline{A}\|_{\mathcal L(\overline{U},Z)}=\|A\|_{\mathcal L(U,Z)}.
\]
\end{proof}

\noindent Application of Theorem 7.1.8 to the Riemann integral:
\[
U=T(I,X),\qquad
\overline{U}=\overline{T(I,X)}=S(I,X)\ \text{w.r.t.}\ \|\cdot\|_\infty
\]
\[
S(I,X)\subset B(I,X)=(Y,\|\cdot\|),\qquad
Z\leftarrow X,\qquad
A=\int_a^b\in\mathcal L\bigl(T(I,X),X\bigr).
\]

\begin{definition}{7.1.9}
The \emph{Riemann integral} is the uniquely defined extension of
$\int_a^b\in\mathcal L\bigl(T(I,X),X\bigr)$ to $\mathcal L\bigl(S(I,X),X\bigr)$ and is denoted by:
\[
\int_a^b f=\int_I f=\int_a^b f(x)\,dx=\int_a^b f\,dx.
\]
\end{definition}

\begin{bemerkung}{7.1.10}
\begin{enumerate}
    \item
    \[
    \left\|\int_a^b f\right\|_X\le(b-a)\|f\|_\infty
    \qquad\forall f\in S(I,X).
    \]

    \item
    \[
    \int_a^b f=\lim_{n\to\infty}\int_a^b f_n,
    \]
    where $(f_n)_{n\in\mathbb N_0}\subset T(I,X)$ is any sequence of step functions with
    \[
    \lim_{n\to\infty}\|f-f_n\|_\infty=0.
    \]
\end{enumerate}
\end{bemerkung}

\begin{satz}{7.1.11 (Approximation of the Riemann integral)}
Let $f\in S(I,Y)$. For every $\varepsilon>0$ there exists a
$\delta(\varepsilon,f)>0$ such that for every partition $Z=(x_0,\ldots,x_n)$ of $I$ with mesh size
\[
\Delta_Z:=\max_{1\le i\le n}(x_i-x_{i-1})<\delta(\varepsilon,f)
\]
and any intermediate points $\xi_j\in[x_{j-1},x_j]$, $1\le j\le n$, we have:
\[
\left\|\int_a^b f-\sum_{j=1}^{n}f(\xi_j)(x_j-x_{j-1})\right\|_Y\le\varepsilon.
\]
\end{satz}

\begin{proof}
\textbf{Step 1:} \textit{Proof for step functions:}

Let $f\in T(I,X)$ and $\hat Z=(\hat x_0,\ldots,\hat x_{\hat n})$ any admissible partition of $I$ for $f$ with
\[
\delta_{\min}:=\min\{(\hat x_i-\hat x_{i-1}):1\le i\le\hat n\}.
\]
Let $\varepsilon>0$ be arbitrary and
\[
\delta=\delta(\varepsilon,f):=\min\left\{\frac{\varepsilon}{4\hat n\|f\|_\infty},\,\delta_{\min}\right\}.
\]
Let $Z=(x_0,\ldots,x_n)$ be any partition of $I$ and $\xi_j\in[x_{j-1},x_j]$, $1\le j\le n$,
any intermediate points with
\[
\Delta_Z:=\max_{1\le j\le n}(x_j-x_{j-1})<\delta.
\]
Since $\Delta_Z<\delta_{\min}$, each interval $[x_{j-1},x_j]$ contains at most one point $\hat x_k$.
We define the refinement $Z':=\hat Z\cup Z=(y_0,\ldots,y_m)$ ($Z'$ is admissible and fine).
Define for $1\le j\le m$ the midpoints $m_j:=(y_{j-1}+y_j)/2$ and
\[
e_j':=f(m_j),\qquad
e_j'':=f(\xi_k)\quad\text{with }k=k(j)\text{ so that }[y_{j-1},y_j]\subset[x_{k-1},x_k]\ni\xi_k.
\]
Since all points where $f$ is discontinuous are contained in $(y_j)_{j=0}^{m}$ and $[x_{k-1},x_k]$ contains
at most one discontinuity, $e_j'\neq e_j''$ can only occur when $y_{j-1}\in\hat Z$ or $y_j\in\hat Z$.
Hence:
\begin{align*}
\#\{j\in\mathbb N_{1,m}:e_j'\neq e_j''\}
&\le\#\bigl(\{j\in\mathbb N_{1,m}:y_{j-1}\in\hat Z\}\cup\{j\in\mathbb N_{1,m}:y_j\in\hat Z\}\bigr)\\
&\le\#\{j\in\mathbb N_{1,m}:y_{j-1}\in\hat Z\}+\#\{j\in\mathbb N_{1,m}:y_j\in\hat Z\}\\
&=\#\hat Z+\#\hat Z=2\hat n.\tag{$*$}
\end{align*}
Furthermore:
\begin{align*}
\int_a^b f-\sum_{k=1}^{n}f(\xi_k)(x_k-x_{k-1})
&=\sum_{j=1}^{m}(y_j-y_{j-1})f(m_j)-\sum_{j=1}^{m}f(\xi_{k(j)})(y_j-y_{j-1})\\
&=\sum_{j=1}^{m}e_j'(y_j-y_{j-1})-\sum_{j=1}^{m}e_j''(y_j-y_{j-1})
=\sum_{j=1}^{m}(e_j'-e_j'')(y_j-y_{j-1}).\tag{$**$}
\end{align*}
By ($*$), at most $2\hat n$ summands in ($**$) are nonzero. Moreover:
\[
\|(y_j-y_{j-1})(e_j'-e_j'')\|_Y
\le|y_j-y_{j-1}|(\|e_j''\|_Y+\|e_j'\|_Y)
<2\|f\|_\infty\delta
\le\frac{\varepsilon}{2\hat n}.
\]
Hence:
\[
\left\|\int_a^b f-\sum_{k=1}^{n}f(\xi_k)(x_k-x_{k-1})\right\|_Y\le\varepsilon.
\]

\medskip
\textbf{Step 2:}
Let now $f\in S(I,X)$ and $\varepsilon>0$. Then there exists $g\in T(I,X)$ with
\[
\|f-g\|_\infty<\frac{\varepsilon}{3(b-a)}.
\]
By Step~1, there exists $\delta=\delta\!\left(\tfrac\varepsilon3,g\right)$ such that for $g$ and $\tfrac\varepsilon3$ the claim of the theorem holds.
Let $Z=(x_0,\ldots,x_n)$ be a partition of $I$ with $\Delta_Z<\delta$ and $\xi_j\in[x_{j-1},x_j]$,
$1\le j\le n$, any intermediate points. Then:
\begin{align*}
&\left\|\int_a^b f-\sum_{j=1}^{n}f(\xi_j)(x_j-x_{j-1})\right\|_Y\\
&\quad\le
\left\|\int_a^b f-\int_a^b g\right\|_Y
+\left\|\int_a^b g-\sum_{j=1}^{n}g(\xi_j)(x_j-x_{j-1})\right\|_Y
+\left\|\sum_{j=1}^{n}\bigl(g(\xi_j)-f(\xi_j)\bigr)(x_j-x_{j-1})\right\|_Y\\
&\quad<\|f-g\|_\infty(b-a)+\frac{\varepsilon}{3}+\sum_{j=1}^{n}\|g-f\|_\infty(x_j-x_{j-1})
<\frac{\varepsilon}{3}+\frac{\varepsilon}{3}+\|g-f\|_\infty(b-a)<\varepsilon.
\end{align*}
\end{proof}

\begin{bemerkung}{7.1.12}
Let $f:I\to X$ be a function. It is called \textbf{Riemann-integrable} if there exists
an $x^*\in X$ such that for every $\varepsilon>0$ there exists a $\delta(\varepsilon,f)>0$ such that for every partition
$Z=(x_0,\ldots,x_n)$ with mesh size $\Delta_Z<\delta$ and any $n$ intermediate points
$\xi_j\in[x_{j-1},x_j]$, $1\le j\le n$, we have:
\[
\left\|x^*-\sum_{i=1}^{n}f(\xi_i)(x_i-x_{i-1})\right\|_X<\varepsilon.
\]
The quantity $x^*$, if it exists, is uniquely determined and is called the Riemann integral of $f$.
Theorem 7.1.11 shows that every jump-continuous function is Riemann-integrable. However, there
exist Riemann-integrable functions that are not jump-continuous. An example is
\[
f:[-1,1]\to\mathbb R,\qquad
f(x)=
\begin{cases}
1 & \text{if } x=\tfrac{1}{n},\ n\in\mathbb N,\\
-1 & \text{if } x=-\tfrac{1}{n},\ n\in\mathbb N,\\
0 & \text{otherwise.}
\end{cases}
\]
The approximating sums
\[
\sum_{i=1}^{n}f(\xi_i)(x_i-x_{i-1})
\]
are also called \emph{Riemann sums}, and one writes symbolically
\[
\int_a^b f=\lim_{\Delta_Z\to 0}\sum_{i=1}^{n}f(\xi_i)(x_i-x_{i-1}).
\]
\end{bemerkung}

\end{document}